\chapter{Relational Database Foundations}
\label{ch:relational-foundations}

\begin{learningobjectives}
After completing this chapter, you will be able to:
\begin{itemize}
  \item define database, schema, table, row, column, and datatype;
  \item distinguish primary, foreign, and composite keys;
  \item reason about one-to-one, one-to-many, and many-to-many relationships;
  \item explain normalization, referential integrity, views, and indexes;
  \item predict how a join can change row counts; and
  \item connect OULAD source grains to the project's raw, staging, and core
    layers.
\end{itemize}
\end{learningobjectives}

\section{The relational model}

A \term{relational database} organizes data into relations commonly presented
as tables. The relational model provides a disciplined way to identify rows,
connect tables, and express transformations \citep{codd1970}. PostgreSQL is the
database system used by this project \citep{postgres16}.

\begin{description}
  \item[Database] A named collection of data and database objects. Here it is
    \texttt{learning\_analytics}.
  \item[Schema] A namespace inside a database. The project uses
    \texttt{raw}, \texttt{staging}, \texttt{core}, \texttt{quality},
    \texttt{analytics}, and \texttt{features}.
  \item[Table] A stored relation made of rows and columns.
  \item[Row] One observation at the table's declared grain.
  \item[Column] One named attribute of every row.
  \item[Datatype] The allowed representation and operations for a column, such
    as \texttt{text}, \texttt{integer}, \texttt{numeric}, or \texttt{boolean}.
\end{description}

\section{Why the project has layers}

\begin{table}[htbp]
\centering
\caption{Database layers and their contracts}
\label{tab:schema-layers}
\begin{tabularx}{\textwidth}{L{0.16\textwidth}L{0.29\textwidth}Y}
\toprule
Schema & Responsibility & Example \\
\midrule
\texttt{raw} & Preserve source columns for loading and reconciliation & Dates and IDs often land as text \\
\texttt{staging} & Normalize names, types, and missing markers & \texttt{NULLIF(date, '?')::integer} \\
\texttt{core} & Enforce stable grains, keys, values, and relationships & Composite primary key for attempts \\
\texttt{quality} & Persist executable checks and severity & Orphan, duplicate, range, and category checks \\
\texttt{analytics} & Define reusable analytical metrics and marts & Weekly engagement and assessment progress \\
\texttt{features} & Construct point-in-time model snapshots & One attempt-week with a future eligible assessment \\
\bottomrule
\end{tabularx}
\end{table}

\begin{plainlanguage}
The layers separate evidence from interpretation. Raw tables answer ``what was
in the file?'' Staging answers ``how should the values be represented?''
Core answers ``what relationships and rules are dependable?'' Later layers
answer analytical questions.
\end{plainlanguage}

\section{Keys}

A \term{primary key} uniquely identifies a stored row and cannot be null. A
\term{foreign key} requires a value to match a key in another table. A
\term{composite key} uses more than one column.

The implemented attempt table declares:

\begin{lstlisting}[language=SQL,caption={Selected key and relationship clauses
from \texttt{sql/03\_core.sql}.}]
CREATE TABLE core.student_attempts (
    code_module text NOT NULL,
    code_presentation text NOT NULL,
    id_student integer NOT NULL,
    ...
    PRIMARY KEY (
        code_module,
        code_presentation,
        id_student
    ),
    FOREIGN KEY (code_module, code_presentation)
        REFERENCES core.module_presentations
        (code_module, code_presentation)
);
\end{lstlisting}

\subsection{Why the key is composite}

A student can appear in multiple modules or presentations.
Therefore, \texttt{id\_student} does not identify an attempt. The three-column
key encodes the actual row meaning.

\begin{validationbox}
To test a proposed key, group by its columns and search for groups with more
than one row:
\begin{lstlisting}[language=SQL]
SELECT
    code_module,
    code_presentation,
    id_student,
    count(*) AS rows_at_key
FROM staging.student_info
GROUP BY 1, 2, 3
HAVING count(*) > 1;
\end{lstlisting}
The saved project run found zero duplicate attempt keys.
\end{validationbox}

\section{Relationships and cardinality}

\term{Cardinality} describes how rows in one table relate to rows in another.

\begin{description}
  \item[One-to-one] Each row matches at most one row on the other side. The
    project expects one registration row for one student attempt.
  \item[One-to-many] One parent can match many children. A module-presentation
    defines multiple assessments.
  \item[Many-to-many] Multiple rows on both sides can match. This relationship
    generally requires an associative table or careful analytical
    pre-aggregation.
\end{description}

An attempt can have many assessments and many VLE interaction rows. If those
two child tables are joined to the attempt simultaneously without
pre-aggregation, each assessment can pair with each interaction. The resulting
row multiplication is mathematically valid as a join result but wrong for most
analytical questions.

\subsection{The five join questions}

Before every join, answer:
\begin{enumerate}
  \item What does one row in the left input represent?
  \item What does one row in the right input represent?
  \item Which columns form the join key?
  \item What should one output row represent?
  \item How should the output row count change?
\end{enumerate}

SQL joins are taught in the forthcoming chapter, but cardinality reasoning must
come first.

\section{Referential integrity}

\term{Referential integrity} means that a child reference points to a valid
parent. For example, every core assessment's module-presentation must exist in
\texttt{core.module\_presentations}.

Not every source relationship can be declared conveniently as a foreign key.
The project supplements core constraints with explicit anti-join checks:

\begin{lstlisting}[language=SQL,caption={Conceptual form of the implemented
orphan-assessment check.}]
SELECT count(*) AS orphan_assessments
FROM staging.assessments AS assessment
LEFT JOIN staging.courses AS course
  USING (code_module, code_presentation)
WHERE course.code_module IS NULL;
\end{lstlisting}

The expected result is zero. A nonzero result means the source contains an
assessment whose parent module-presentation is absent.

\section{Nulls and constraints}

SQL \texttt{NULL} means the value is unknown or absent under the database
representation. It is not zero and not an empty string. Conditions such as
\texttt{score = NULL} do not work; use \texttt{score IS NULL}.

The core layer adds \term{constraints}:
\begin{itemize}
  \item \texttt{NOT NULL} for required values;
  \item \texttt{CHECK (score BETWEEN 0 AND 100)};
  \item accepted assessment types;
  \item accepted final-result categories; and
  \item positive presentation lengths and studied credits.
\end{itemize}

Constraints prevent invalid future writes. Quality queries also document and
profile the source before or beyond what constraints can express.

\section{Normalization and denormalization}

\term{Normalization} separates facts so that each fact is stored at an
appropriate grain. Module-presentation details, assessments, attempts,
resources, submissions, and interactions therefore live in different core
tables.

\term{Denormalization} deliberately combines or aggregates facts for repeated
analytical use. The project's materialized weekly engagement mart is
denormalized to one active attempt-week. It is not a replacement for the core
tables; it is a derived analytical product.

\section{Views, materialized views, and indexes}

\begin{description}
  \item[View] A saved query evaluated when read. The staging layer uses views so
    source-faithful raw tables remain available.
  \item[Materialized view] A saved query whose result rows are stored and later
    refreshed or rebuilt. The project uses these for expensive analytical
    marts.
  \item[Index] An auxiliary structure that can speed selected access patterns
    while using storage and adding write cost.
\end{description}

These concepts are fully developed in forthcoming SQL chapters. At this stage,
remember that an index does not change the meaning or grain of a table.

\section{Knowledge check}

\begin{enumerate}
  \item Why does the project preserve raw text before casting values?
  \item What rule does the attempt composite key enforce?
  \item Give one one-to-many OULAD relationship.
  \item Explain how joining two child tables can multiply rows.
  \item What is the difference between a view and a materialized view?
  \item Why are both constraints and quality queries useful?
\end{enumerate}

\begin{exercisebox}{Exercise 4.1: Diagnose a dangerous join plan}
A learner proposes joining attempts, assessments, and VLE interactions in one
query using only module and presentation. Before writing SQL:
\begin{enumerate}
  \item state each input grain;
  \item describe the expected multiplication;
  \item identify the missing student key where applicable;
  \item propose which child data should be pre-aggregated; and
  \item define the desired output grain.
\end{enumerate}
\solutionref{sol:dangerous-join}
\end{exercisebox}

\begin{exercisebox}{Exercise 4.2: Classify database rules}
For each rule, choose primary key, foreign key, check constraint, null rule, or
quality query, and justify the choice:
\begin{enumerate}
  \item an assessment weight must be between 0 and 100;
  \item an interaction site should exist in the same presentation;
  \item an attempt is unique by module, presentation, and student;
  \item missing deprivation band should remain visible; and
  \item the official archive should contain 22 course rows.
\end{enumerate}
\solutionref{sol:database-rules}
\end{exercisebox}

\begin{commonmistakes}
\begin{itemize}
  \item Declaring a key before defining row grain.
  \item Believing a successful join preserves the intended row count.
  \item Replacing unknown values with zero.
  \item Treating indexes as data-cleaning tools.
  \item Removing the raw layer after typed tables are built.
\end{itemize}
\end{commonmistakes}

\transferheading
Layered relational design applies beyond OULAD. Preserve a source-faithful
landing layer, make type and missingness decisions visible, enforce dependable
core contracts, and build denormalized analytical products only after the
grain is stable.

\begin{summarybox}
The project uses six PostgreSQL schemas to separate source evidence, typed
representations, constrained core data, validation, analytics, and features.
Keys express row grain; foreign keys and quality checks express relationships.
Nulls, constraints, cardinality, and row multiplication must be reasoned about
before analytical SQL is trusted.
\end{summarybox}
